\chapter{Quantum Channels, Distances, and Statistical Estimates}
\label{ch:lec02}

\begin{objectives}
\begin{itemize}[leftmargin=1.4em,nosep]
  \item Describe quantum operations as completely positive trace-preserving maps.
  \item Compare states and channels using trace distance, fidelity, and the diamond norm.
  \item Apply concentration inequalities to bound measurement error.
  \item Use statistical amplification and the Cram\'er--Rao bound to reason about shot budgets.
\end{itemize}
\end{objectives}

\todo{Write the introductory paragraph orienting the lecture.}

\section{Quantum channels (CPTP maps)}
\label{sec:lec02-1}
\todo{Draft this section.}

\section{Distances between states and channels}
\label{sec:lec02-2}
\todo{Draft this section.}

\section{Elementary concentration inequalities}
\label{sec:lec02-3}
\todo{Draft this section.}

\section{Statistical amplification and shot budgets}
\label{sec:lec02-4}
\todo{Draft this section.}

\section{Parameter estimation and the Cram\'er--Rao bound}
\label{sec:lec02-5}
\todo{Draft this section.}

\exercisesection
\begin{exercise}
\todo{Add exercise.}
\end{exercise}
\begin{exercise}
\todo{Add exercise.}
\end{exercise}
\begin{exercise}
\todo{Add exercise.}
\end{exercise}

\readingsection
\begin{itemize}[leftmargin=1.4em]
  \item Lin Lin, \emph{live notes}, Ch.~3 (probability, channels, distances).
  \item Lin Lin, \emph{live notes}, Ch.~8 (statistical estimates).
\end{itemize}
